\chapter{LLMs for Credit Risk Analysis}
\label{ch:credit-risk}

\begin{remark}[Learning Objectives]
After completing this chapter, the reader should be able to:
\begin{enumerate}
  \item Describe the main sources of credit data, identify the regulatory constraints
        that govern their use (GDPR, FCRA, ECOA), and apply standard preprocessing
        pipelines to handle class imbalance and missing values.
  \item Formulate default prediction as a text classification task and fine-tune a
        pre-trained language model on credit-related data using parameter-efficient methods.
  \item Extract implied default probabilities from a language model using structured
        generation and calibrate those probabilities with isotonic regression and
        Platt scaling.
  \item Model complex household decisions under uncertainty using LLMs as
        decision-support agents grounded in utility theory and bounded rationality.
  \item Design and parameterise persona agents that simulate heterogeneous household
        financial behaviour, and calibrate their outputs against observed survey data.
  \item Evaluate a credit risk model using standard discrimination metrics (AUROC,
        KS statistic, Gini coefficient) and validate it within the SR~11-7 model
        risk management framework.
  \item Apply SHAP-based interpretability methods to an LLM credit decision system
        and produce a compliant explanation consistent with GDPR Article~22.
  \item Architect and deploy a production credit risk application with monitoring,
        drift detection, human-in-the-loop review, and regulatory documentation.
\end{enumerate}
\end{remark}

% ---------------------------------------------------------------
\section{Credit Data: Sources, Privacy, and Regulatory Constraints}
\label{sec:credit-data}

[Placeholder]

\subsection{Credit Bureau Data and Alternative Data Sources}
\label{subsec:credit-bureau}

[Placeholder]

\subsection{Regulatory Constraints on Credit AI: GDPR, FCRA, and ECOA}
\label{subsec:credit-regulation}

[Placeholder]

\subsection{Preprocessing: Anonymisation, Class Imbalance, and Missing Data}
\label{subsec:credit-preprocessing}

[Placeholder]

% ---------------------------------------------------------------
\section{Credit Risk Modelling with LLMs}
\label{sec:credit-modelling}

[Placeholder]

\subsection{From Logistic Regression to Language Models: The Credit Risk Arc}
\label{subsec:credit-arc}

[Placeholder]

\subsection{Fine-Tuning for Default Prediction: Task Formulation and Training}
\label{subsec:credit-finetuning}

[Placeholder]

\subsection{Structured Generation and Implied Default Probabilities}
\label{subsec:structured-gen}

[Placeholder]

\subsection{Probability Calibration: Isotonic Regression and Platt Scaling}
\label{subsec:calibration}

[Placeholder]

% ---------------------------------------------------------------
\section{Household Decisions Under Uncertainty}
\label{sec:household-decisions}

[Placeholder]

\subsection{Bounded Rationality and the Household Decision Problem}
\label{subsec:bounded-rationality}

[Placeholder]

\subsection{LLMs as Decision-Support Agents for Complex Financial Choices}
\label{subsec:decision-support}

[Placeholder]

\subsection{Utility-Theoretic Framing and Preference Elicitation}
\label{subsec:utility-elicitation}

[Placeholder]

% ---------------------------------------------------------------
\section{Persona Agents and Behavioural Simulation}
\label{sec:persona-agents}

[Placeholder]

\subsection{Defining Agent Personas: Income, Education, Risk Aversion, and Financial Literacy}
\label{subsec:persona-design}

[Placeholder]

\subsection{Multi-Persona Dialogue and Population-Level Disagreement}
\label{subsec:persona-dialogue}

[Placeholder]

\subsection{Applications: Loan Uptake, Mortgage Choice, and Retirement Planning}
\label{subsec:persona-apps}

[Placeholder]

\subsection{Calibrating Simulated Behaviour Against Survey Data}
\label{subsec:persona-calibration}

[Placeholder]

% ---------------------------------------------------------------
\section{Evaluation and Model Risk Management}
\label{sec:model-risk}

[Placeholder]

\subsection{Credit-Specific Metrics: AUROC, KS Statistic, and Gini Coefficient}
\label{subsec:credit-metrics}

[Placeholder]

\subsection{SR~11-7 Model Risk Framework Applied to LLM Governance}
\label{subsec:sr117}

[Placeholder]

\subsection{Backtesting and Through-the-Cycle Validation}
\label{subsec:backtesting}

[Placeholder]

\subsection{Interpretability: SHAP for LLM-Based Credit Decisions}
\label{subsec:shap-credit}

[Placeholder]

% ---------------------------------------------------------------
\section{Deployment of a Credit Risk Application}
\label{sec:credit-deployment}

[Placeholder]

\subsection{System Architecture: API Design, Database Integration, and Inference Serving}
\label{subsec:system-arch}

[Placeholder]

\subsection{Monitoring: Data Drift, Model Degradation, and Alerting}
\label{subsec:monitoring}

[Placeholder]

\subsection{Human-in-the-Loop Review Workflows}
\label{subsec:hitl-credit}

[Placeholder]

\subsection{Regulatory Submission and Model Documentation}
\label{subsec:model-doc}

[Placeholder]
